\clearpage
\SourceLeaf{228}\pdfbookmark[1]{Sectio secunda.}{liber-secundus-sectio-secunda}
\begin{center}
{\large SECTIO SECUNDA.}\par\medskip
\emph{Determinatio orbitae e quatuor observationibus, quarum duae tantum completae sunt.}\par\bigskip
\hypertarget{native-228-article164}{}\pdfbookmark[2]{164.}{article-164}164.
\end{center}
Iam in ipso limine Libri secundi (art. 115) declaravimus, usum problematis
in Sect. praec. pertractati ad eas orbitas limitari, quarum inclinatio nec
evanescit, nec nimis exigua est, determinationemque orbitarum parum inclinatarum
necessario quatuor observationibus superstrui debere. Quatuor autem
observationes completae, quum octo aequationibus aequivaleant, incognitarumque
numerus ad sex tantum ascendat, problema plus quam determinatum redderent:
quapropter a duabus observationibus latitudines (sive declinationes) seponere
oportebit, ut datis reliquis exacte satisfieri possit. Sic oritur problema,
cui haec Sectio dicata erit: solutio autem, quam hic trademus, non solum ad
orbitas parum inclinatas patebit, sed etiam ad orbitas inclinationis quantumvis
magnae pari successu applicari poterit. Etiam hic, perinde ut in problemate
Sect. praec., casum eum, ubi orbitae dimensiones approximatae iam in potestate
sunt, segregare oportet a determinatione prima orbitae penitus adhuc incognitae:
ab illo initium faciemus.
